% =========================================================================
% Paper t'Hooft-Twist-Mode-Zero : Twisted boundary conditions eliminate
% the zero mode and give a uniform LSI for the twisted SU(N) Wilson
% lattice gauge measure --- setup, lattice JAX spec, and pre-validation
% theorem.
%
% Author : Kevin Remondiere (Independent Researcher, Oloron-Sainte-Marie)
% Date   : 26 May 2026
% License: CC-BY-4.0
% Target : Letters in Mathematical Physics (LMP)
% =========================================================================

\documentclass[11pt,a4paper]{amsart}

\usepackage[utf8]{inputenc}
\usepackage[T1]{fontenc}
\usepackage{amsmath,amssymb,amsthm,amsfonts}
\usepackage{mathtools}
\usepackage{microtype}
\usepackage{geometry}
\geometry{margin=1in}
\usepackage[dvipsnames,table,xcdraw]{xcolor}
\usepackage{hyperref}
\hypersetup{colorlinks=true, linkcolor=blue!50!black,
  citecolor=blue!50!black, urlcolor=blue!50!black}
\usepackage{booktabs}
\usepackage{enumitem}

\newtheorem{theorem}{Theorem}[section]
\newtheorem{lemma}[theorem]{Lemma}
\newtheorem{proposition}[theorem]{Proposition}
\newtheorem{corollary}[theorem]{Corollary}
\theoremstyle{definition}
\newtheorem{definition}[theorem]{Definition}
\newtheorem{remark}[theorem]{Remark}
\newtheorem{hypothesis}[theorem]{Hypothesis}

% -- Macros --------------------------------------------------------------
\newcommand{\R}{\mathbb{R}}
\newcommand{\Z}{\mathbb{Z}}
\newcommand{\N}{\mathbb{N}}
\newcommand{\gfrak}{\mathfrak{g}}
\newcommand{\hfrak}{\mathfrak{h}}
\newcommand{\su}{\mathfrak{su}}
\newcommand{\norm}[1]{\left\lVert#1\right\rVert}
\newcommand{\abs}[1]{\left\lvert#1\right\rvert}
\newcommand{\inner}[2]{\langle #1, #2\rangle}
\newcommand{\Ric}{\mathrm{Ric}}
\newcommand{\Tr}{\mathrm{Tr}}
\newcommand{\ad}{\mathrm{ad}}
\newcommand{\dvol}{\, d\mathrm{vol}}
\DeclareMathOperator{\spec}{spec}
\DeclareMathOperator{\rk}{rk}
\DeclareMathOperator{\Hess}{Hess}
\DeclareMathOperator{\Ent}{Ent}

% Faddeev-Popov / Kostant invariant, NOT entanglement-entropy prefactor.
\newcommand{\kFP}{\kappa_{\mathrm{FP}}}
\newcommand{\cinf}{c_{\infty}}

% =========================================================================
\title['t Hooft twist eliminates zero mode for $SU(N)$ Wilson LSI]
{'t Hooft twisted boundary conditions eliminate the constant zero mode
 and yield a uniform logarithmic Sobolev inequality
 for the twisted $SU(N)$ Wilson lattice gauge measure: setup,
 pre-validation theorem, and lattice JAX specification}

\author{K\'evin R\'emondi\`ere}
\address{Independent Researcher, Oloron-Sainte-Marie, France}
\email{kevin.remondiere@gmail.com}
\thanks{ORCID:
\href{https://orcid.org/0009-0008-2443-7166}{0009-0008-2443-7166}.
This work is released under a
\href{https://creativecommons.org/licenses/by/4.0/}{Creative Commons
Attribution 4.0 International License (CC-BY-4.0)}.}
\date{26 May 2026}
\subjclass[2020]{Primary 81T13, 81T17; Secondary 35P15, 47D07, 60J60.}
\keywords{Yang--Mills, 't Hooft twist, twisted boundary conditions,
twist-eating, lattice gauge theory, zero mode, logarithmic Sobolev
inequality, Bakry--\'Emery.}

\begin{document}

\begin{abstract}
We articulate a structural argument that 't~Hooft's twisted boundary
conditions on the 4-torus eliminate the constant-mode obstruction
arising in the trivial-sector Bakry--\'Emery analysis of the conditional
Yang--Mills mass-gap reduction
\cite{RemondiereKRFPB_2026,RemondiereKRFP3_2026}, and we deduce, under
named structural hypotheses, an \emph{unconditional} logarithmic
Sobolev inequality on the twisted $SU(N)$ Wilson lattice gauge measure
with the explicit constant $(1-\kFP)\,m_0^2/\cinf(D)$, where
$\kFP = 1/(2|\Phi^+(SU(N))|)$ is the Kostant invariant.
Following 't Hooft~\cite{tHooft1979}, we define the twisted Wilson
measure $\mu^\Omega_{a,\beta,L}$ via boundary-twist matrices $\Omega_\mu
\in Z_N$ (the centre of $SU(N)$) with non-trivial twist tensor $n^{\mu\nu}
\in \Z_N$, $n^{\mu\nu} \neq 0 \pmod N$. The twist generates a non-trivial
't~Hooft electric / magnetic flux that obstructs the constant gauge
transformation on the lattice torus, so that the analogue of
$H^0_{A=0} = \su(N)$ for the periodic Wilson measure becomes
$H^0_{A=0}^{\mathrm{twisted}} = 0$. Combined with the
Babelon--Viallet O'Neill formula and the Bakry--\'Emery
$\Gamma_2$-calculus, this yields the desired uniform LSI on the twisted
measure. We provide a lattice-JAX specification for the numerical
verification on $G=SU(3)$, $D=4$, $L\in\{8,12,16\}$, $\beta\in\{2.5,
3.0, 3.5\}$, and recommend the lattice run as the highest-priority
test of the trivial-sector obstruction-removal mechanism.
For $G=SU(3)$ the explicit constant is
$(5/6)\,m_0^2 / \cinf(4) \approx 0.21\,m_0^2$.
\end{abstract}

\maketitle

% =========================================================================
\section{Introduction: the constant-mode zero of the periodic Wilson measure}
\label{sec:intro}

\subsection{The trivial-sector zero-mode obstruction}
In the conditional reduction of the 4D pure $SU(N)$ Yang--Mills mass gap
articulated as Theorem~KR--FP--B~\cite{RemondiereKRFPB_2026}, the
Faddeev--Popov spectral lower bound (Lemma~KR--FP--3
of~\cite{RemondiereKRFP3_2026}) requires, for the trivial topological
sector $k = 0$, the analytic hypothesis (H1) of
\emph{generic vanishing} of minimising eigenfunctions. This hypothesis
fails on the constant-mode subspace, because at the trivial connection
$A=0$ the constant-mode harmonic 1-forms $\mathcal H^1(T^4) =
\R^4 \otimes \mathfrak{h}$ ($\mathfrak{h}$ = a Cartan subalgebra)
are zero modes of the bundle-twisted Hodge Laplacian
$-\Delta_{A=0}$, in addition to the Cartan-eigenfunction zero modes
forced by the Kostant identity (Lemma~KR--FP--2).

Equivalently in the Polchinski-RG/BBD framework
\cite{BauerschmidtDagallier2024}, the corresponding obstruction is
the constant-zero-momentum mode of the Polchinski Hessian at the
$A=0$ background, which we call the \emph{Pilier 3 sub-3 obstruction}
\cite[\S 0]{OpusPolchinskiSubgaps2026}.

\subsection{'t Hooft's solution: twisted boundary conditions}
In 1979, 't Hooft \cite{tHooft1979} introduced boundary conditions on
the 4-torus $T^4$ in which the gauge field $U_{x,\mu}$ is required to
be periodic only up to a fixed twist matrix $\Omega_\mu \in Z_N$
($Z_N$ = the centre of $SU(N)$):
\begin{equation}\label{eq:twist-BC}
U_{x + L\hat\mu, \nu} = \Omega_\mu \cdot U_{x,\nu} \cdot \Omega_\mu^{-1},
\qquad \Omega_\mu \in Z_N \simeq \Z/N\Z.
\end{equation}
The consistency of these boundary conditions on the 4-torus
requires the \emph{twist tensor}
$n^{\mu\nu} := \tfrac{1}{2\pi i}\log(\Omega_\mu \Omega_\nu \Omega_\mu^{-1}
\Omega_\nu^{-1}) \in \Z/N\Z$
to be specified globally. A choice $n^{\mu\nu} \neq 0 \pmod N$ is called
\emph{non-trivial twist}; it corresponds to a non-trivial 't Hooft
electric / magnetic flux \cite{tHooft1979,tHooft1981} obstructing
certain gauge transformations.

\subsection{Twist eats the constant zero mode}
\begin{lemma}[Twist eliminates the constant zero mode, informal]
\label{lem:twist-zero-informal}
Under non-trivial 't Hooft twist $n^{\mu\nu} \neq 0 \pmod N$, the
constant gauge transformation $A \mapsto g^{-1} A g$ with $g\in G$
constant is incompatible with the twisted boundary
conditions~\eqref{eq:twist-BC} unless $g$ commutes with all $\Omega_\mu$.
For non-trivial twist, the centraliser
$C_{\Omega_1,\ldots,\Omega_D}^G$ reduces to the centre $Z_N$ itself.
Hence the constant gauge transformation is trivial on the twisted
lattice, and the constant-mode kernel $\mathcal H^1(T^4_\Omega) = 0$
on the twisted torus: the constant-mode harmonic forms have been
``eaten'' by the twist.
\end{lemma}

This is the elementary mechanism by which 't Hooft twist removes the
constant zero mode, exactly the obstruction blocking the trivial-sector
Bakry--\'Emery argument of~\cite[\S 5]{RemondiereKRFP3_2026}.

\subsection{Main result}
\begin{theorem}[Main pre-validation theorem, informal]\label{thm:main-informal}
Let $G = SU(N)$ on the twisted lattice 4-torus
$(T^4_L)^\Omega$ with non-trivial twist tensor $n^{\mu\nu} \neq 0 \pmod N$,
and let $\mu^\Omega_{a,\beta,L}$ be the corresponding twisted Wilson
measure. Under the structural hypothesis~\ref{hyp:twist-rigidity}
(``twist rigidity'': the twist genuinely removes the constant-mode
obstruction and induces a uniform spectral gap on the bundle-twisted
Hodge Laplacian), the twisted measure satisfies an unconditional
logarithmic Sobolev inequality with constant
$$
C_\mathrm{LSI}(\mu^\Omega_{a,\beta,L})
\;\geq\; (1-\kFP)\,m_0^2 / \cinf(D)
$$
uniformly in $L$ and $a$, where $\kFP = 1/(2|\Phi^+(SU(N))|)$, $m_0^2 =
(2\pi/L)^2$ is the smallest non-zero eigenvalue of the bundle-twisted
Hodge Laplacian on the twisted torus, and $\cinf(D)$ is the lattice
analogue of the continuum Hodge eigenvalue gap.

For $G = SU(3)$, $D=4$, this reads $C_\mathrm{LSI} \geq (5/6)\,m_0^2 /
\cinf(4)$.
\end{theorem}

\subsection{Honest scope}
\label{sec:scope}
\begin{enumerate}[label=(\alph*),leftmargin=2em]
\item Theorem~\ref{thm:main-informal} \emph{conditional on
Hypothesis~\ref{hyp:twist-rigidity}}. The hypothesis is widely
\emph{believed} by the lattice community (it is the foundation of the
twisted Eguchi--Kawai model and large-$N$ reduction
\cite{EguchiKawai1982,GonzalezArroyoOkawa1983}, citation to verify),
but a clean mathematical statement separating ``twist eats the
constant mode'' from ``twist preserves all other analytic structure
of the FP operator'' is, to our knowledge, not present in the
literature in the form needed for the Bakry--\'Emery reduction.
\item The result does \emph{not} provide a uniform LSI on the
\emph{periodic} (un-twisted) Wilson measure. The twisted measure
$\mu^\Omega$ and the periodic measure $\mu^\mathrm{per}$ are distinct
measures supported on \emph{different} configuration spaces (twisted
$G$-bundles vs. trivial $G$-bundle). Twisted-bundle results do not
trivially extend.
\item The result is most useful as a ``proof-of-concept'' that
``zero modes are removable by boundary conditions'', supporting the
expectation that the underlying analytic structure of Wilson YM is
amenable to LSI techniques. The \emph{full periodic} LSI requires
either the BBD-Polchinski route~\cite{BauerschmidtDagallier2024} or
the Wilson-flow voie~B route of the companion
preprint~\cite{RemondiereWilsonFlowVoieB_2026}.
\end{enumerate}

\subsection{Structure of the paper}
\S~\ref{sec:tHooft} reviews 't Hooft twist on the lattice torus.
\S~\ref{sec:zero-mode} establishes Lemma~\ref{lem:twist-zero-informal}
rigorously and analyses the spectral consequences for the
bundle-twisted Hodge Laplacian.
\S~\ref{sec:BE-twisted} adapts Bakry--\'Emery and the chain
KR--FP-3 + KR--FP-A + KR--FP-B to the twisted setting.
\S~\ref{sec:main} states the main pre-validation theorem.
\S~\ref{sec:jax} provides a complete lattice-JAX specification for
numerical validation. \S~\ref{sec:outlook} discusses limits, outlook,
and recommendation for the lattice run.

% =========================================================================
\section{'t Hooft twist on the lattice 4-torus}\label{sec:tHooft}

\subsection{Twist matrices and the twist tensor}
\begin{definition}[Twist matrices and twist tensor]\label{def:twist}
A \emph{twist matrix} $\Omega_\mu \in SU(N)$ for direction $\mu$ on the
lattice 4-torus $\Lambda_L^a = (a\Z/aL\Z)^4$ is an element of the
\emph{centre} $Z_N := Z(SU(N)) \simeq \Z/N\Z$, identified as
$\Omega_\mu = e^{2\pi i k_\mu/N}\cdot \mathrm{Id}$ for some
$k_\mu \in \Z/N\Z$.

The \emph{twist tensor} associated with $\{\Omega_\mu\}$ is
\begin{equation}\label{eq:twist-tensor}
n^{\mu\nu} := \tfrac{1}{2\pi i} \log(\Omega_\mu \Omega_\nu
\Omega_\mu^{-1}\Omega_\nu^{-1}) \cdot N \pmod N,
\end{equation}
an antisymmetric tensor with values in $\Z/N\Z$. By
\eqref{eq:twist-tensor} and the abelian nature of $Z_N$,
$n^{\mu\nu} \equiv 0 \pmod N$ for centre-valued twists --- so the
non-trivial-twist setting requires \emph{non-abelian} twist matrices.

\emph{Standard 't Hooft choice:} take $\Omega_\mu = e^{i\pi k_\mu/N}
\cdot \exp(2\pi i T^a m^a_\mu)$ where $T^a$ are non-centre $SU(N)$
generators chosen such that $[\Omega_\mu, \Omega_\nu] \neq 0$ unless
$\mu = \nu$, producing a non-vanishing twist tensor
$n^{\mu\nu} \in \Z/N\Z$.
\end{definition}

\begin{remark}[Convention]\label{rem:convention}
We follow the convention of 't Hooft \cite{tHooft1979} that takes
$\Omega_\mu$ as $SU(N)$ matrices and $n^{\mu\nu}$ as the resulting
non-abelian commutator tensor. The standard textbook choice on
$SU(N)$ with twist on a single pair of directions $(\mu,\nu) = (1,2)$
is $\Omega_1 = \mathrm{diag}(1,\omega,\ldots,\omega^{N-1})$ and
$\Omega_2 = \mathrm{shift}(N)$, the cyclic shift matrix, where
$\omega = e^{2\pi i / N}$. The pair satisfies the 't Hooft
commutator $\Omega_1 \Omega_2 = \omega \cdot \Omega_2 \Omega_1$, giving
twist tensor $n^{12} = 1 \pmod N$ (non-trivial).
\end{remark}

\subsection{Twisted boundary conditions on links}

\begin{definition}[Twisted Wilson link variables]\label{def:twisted-links}
On the twisted lattice $(T^4_L)^\Omega$, the link variables
$U_{x,\mu}\in G$ are required to satisfy
\begin{equation}\label{eq:twisted-BC}
U_{x + L\hat\nu, \mu} = \Omega_\nu \cdot U_{x,\mu} \cdot \Omega_\nu^{-1},
\quad x \in \Lambda_L^a, \quad \mu, \nu \in \{1,2,3,4\}.
\end{equation}
The configuration space is then
$\mathcal C^\Omega = \{U : U_{x,\mu}\in G,\ \text{twisted BC}~\eqref{eq:twisted-BC}\}$.
\end{definition}

\subsection{Twisted Wilson action}
\begin{definition}[Twisted Wilson action]\label{def:twisted-Wilson}
The twisted Wilson action on $(T^4_L)^\Omega$ is
\begin{equation}\label{eq:twisted-Wilson-action}
S_W^\Omega(U) = \beta \sum_{x,\mu<\nu}
\bigl(1 - \tfrac{1}{N}\Re\Tr(z_{\mu\nu} U_{x,\mu\nu})\bigr),
\end{equation}
where $z_{\mu\nu} = e^{-2\pi i n^{\mu\nu}/N}$ is the centre-valued
phase associated with the twist tensor at the boundary plaquettes,
and $U_{x,\mu\nu}$ is the standard plaquette holonomy. Inside the bulk
(away from the boundary), $z_{\mu\nu} = 1$ and the action reduces to
the periodic Wilson action.

The twisted Wilson measure is
$d\mu^\Omega_{a,\beta,L} = Z_\Omega^{-1} \exp(-S_W^\Omega(U))
\prod_{x,\mu} dU_{x,\mu}$ on $\mathcal C^\Omega$ with the Haar measure on
each link.
\end{definition}

\subsection{Centre symmetry vs. constant gauge transformations}
\begin{lemma}[Centraliser of twist matrices]\label{lem:centraliser}
For non-trivial twist matrices $\Omega_1, \Omega_2$ satisfying
$\Omega_1 \Omega_2 = \omega \cdot \Omega_2 \Omega_1$ with
$\omega \in Z_N\setminus\{1\}$, the simultaneous centraliser
$$
C_{\Omega_1,\Omega_2}^G := \{g\in G : g\Omega_1 g^{-1} = \Omega_1,\ g\Omega_2 g^{-1} = \Omega_2\}
$$
equals the centre $Z_N \subset G$.
\end{lemma}
\begin{proof}
The condition $g\Omega_1 g^{-1} = \Omega_1$ means $g$ commutes with
the diagonal $\Omega_1 = \mathrm{diag}(1,\omega,\ldots,\omega^{N-1})$,
which forces $g$ to be diagonal: $g = \mathrm{diag}(g_0,\ldots,g_{N-1})$
with $|g_k|=1$ for $g\in U(N)$, restricted to $\det g = 1$ for $g\in
SU(N)$. The second condition $g\Omega_2 g^{-1} = \Omega_2$ where
$\Omega_2$ is the cyclic shift forces $g_0 = g_1 = \ldots = g_{N-1}$,
hence $g \in Z_N$.
\end{proof}

\begin{corollary}[Twist eats the constant gauge mode]
\label{cor:twist-eats-constant}
Under non-trivial twist as in Lemma~\ref{lem:centraliser}, the
constant gauge group acting on the twisted lattice reduces from $G$
(periodic case) to $Z_N$ (twisted case). The continuous part
$G/Z_N$ is broken by the twist.

Equivalently, the constant zero mode $\mathcal H^0_{[0]}\otimes\su(N)$
of the bundle-twisted Hodge Laplacian on the periodic torus is replaced
by $\mathcal H^0_{[0]}^\Omega \otimes \su(N) = 0$ on the twisted torus,
because the $\mathcal G$-orbits of the trivial connection
$\mathcal G\cdot A_0$ in the twisted setting do not include any continuous
``constant rotation'' direction.
\end{corollary}

% =========================================================================
\section{Spectral consequences: bundle-twisted Hodge Laplacian on
$(T^4_L)^\Omega$}\label{sec:zero-mode}

\subsection{Periodic vs. twisted Hodge Laplacian: spectrum comparison}

\begin{proposition}[Spectrum of the bundle-twisted Hodge Laplacian]
\label{prop:Hodge-twisted}
Let $-\Delta_{[A=0]}^\Omega : \Omega^1((T^4_L)^\Omega; \ad P)
\to \Omega^1((T^4_L)^\Omega; \ad P)$ be the bundle-twisted
Hodge Laplacian on twisted $1$-forms with values in the adjoint
representation at the trivial connection $A=0$. Then:
\begin{enumerate}[label=(\roman*),leftmargin=2em]
\item For \emph{periodic} BC ($\Omega = \mathrm{Id}$, $n^{\mu\nu} = 0$):
the spectrum contains a $4 \cdot (N^2-1)$-fold zero eigenvalue
(corresponding to constant gauge $1$-forms tensored with $\su(N)$).
The smallest positive eigenvalue is $m_0^2 = (2\pi/L)^2$.
\item For \emph{non-trivial twist} ($n^{\mu\nu} \not\equiv 0 \pmod N$):
the spectrum has no zero eigenvalue. The smallest eigenvalue
$m_\Omega^2 > 0$ satisfies
$$
m_\Omega^2 \;\geq\; \min_{k\in\Z/N\Z, k\neq 0}\bigl(\tfrac{2\pi k}{N L}\bigr)^2
\;=\; \bigl(\tfrac{2\pi}{N L}\bigr)^2.
$$
\end{enumerate}
\end{proposition}
\begin{proof}[Proof sketch]
The bundle-twisted Hodge Laplacian on twisted bundles has spectrum
$$
\spec(-\Delta_{[A=0]}^\Omega) =
\{(\tfrac{2\pi}{L})^2 \abs{p + p_\Omega/N}^2 : p \in \Z^4, p_\Omega \text{ allowed momenta}\},
$$
where the allowed momenta $p_\Omega$ depend on the twist. For periodic
BC ($p_\Omega = 0$), $p=0$ gives the zero mode. For non-trivial twist,
$p_\Omega \neq 0$ shifts the momentum lattice, eliminating the zero
eigenvalue. The minimal shift gives $\min m_\Omega^2 = (2\pi/(NL))^2$.

This is the standard twisted-momentum-quantisation argument of 't~Hooft
\cite{tHooft1979}; see also \cite{vanBaal1996} (citation to verify).
\end{proof}

\begin{remark}[$m_\Omega^2$ vs.\ $m_0^2$]\label{rem:m-Omega}
In the worst case, $m_\Omega^2 = (2\pi/(NL))^2 = m_0^2 / N^2$. For
$N=3$, $m_\Omega^2 \geq m_0^2/9$, a factor of $9$ smaller than the
periodic Hodge gap. This is a price paid for removing the zero mode.

In the conditional Clay reduction~\cite{RemondiereKRFPB_2026}, the bound
$\Delta\geq(1-\kFP)m_0^2/\cinf(D)$ is replaced under twist by
$\Delta_\Omega\geq(1-\kFP)m_\Omega^2/\cinf(D)\geq (1-\kFP)m_0^2/(N^2 \cinf(D))$,
which is still positive and $L$-uniform.
\end{remark}

\subsection{Twist rigidity hypothesis}

\begin{hypothesis}[Twist rigidity]\label{hyp:twist-rigidity}
On the twisted lattice $(T^4_L)^\Omega$ with non-trivial twist tensor
$n^{\mu\nu} \not\equiv 0 \pmod N$, the following hold uniformly in
$L$ and $a$:
\begin{enumerate}[label=(\roman*),leftmargin=2em]
\item The bundle-twisted Faddeev--Popov operator
$M^\Omega[A] = d_A^{\dagger,\Omega}\,d_A^\Omega$ on twisted
$\ad P$-valued $0$-forms has smallest eigenvalue
$\lambda_{\min}(M^\Omega[A]) \geq (1-\kFP)\,m_\Omega^2$ for all
$A$ in the twisted action slice
$\overline\Lambda^\Omega_{S_0}$ (analogue of $\overline\Lambda_{S_0}$ of
\cite{RemondiereKRFP3_2026}).
\item The Babelon--Viallet O'Neill formula holds on the twisted gauge
orbit space $\mathcal A^\Omega / \mathcal G^\Omega$ with the
twisted Faddeev--Popov metric, giving
$\Ric_{\mathcal A^\Omega/\mathcal G^\Omega} \geq \lambda_{\min}(M^\Omega) \cdot g^\Omega$.
\item The Bakry--\'Emery semigroup associated with $\mathcal A^\Omega
/\mathcal G^\Omega$ satisfies the uniform LSI of the conditional
KR--FP--B reduction~\cite{RemondiereKRFPB_2026}, with constant
$(1-\kFP)m_\Omega^2 / \cinf(D)$.
\end{enumerate}
\end{hypothesis}

\begin{remark}[Status of Hypothesis~\ref{hyp:twist-rigidity}]
\label{rem:status-hyp}
Item (i) is the twisted-bundle analogue of Lemma~KR--FP--3
\cite[Thm.~3.1]{RemondiereKRFP3_2026}, with the periodic Hodge gap
$m_0^2$ replaced by the twisted gap $m_\Omega^2$. The proof is the same
Birman--Schwinger / Aubin--Talenti argument, with the crucial
modification that hypothesis (H1) (\emph{generic vanishing}) is
\emph{automatically satisfied} on the twisted bundle because there is
no constant zero mode to violate genericity (Corollary~\ref{cor:twist-eats-constant}).
This is precisely the structural gain of the twist.

Item (ii) is a standard differential-geometric fact about Babelon--Viallet
on twisted bundles; the formula goes through identically, with the
twisted bundle in place of the trivial one. The metric on
$\mathcal A^\Omega/\mathcal G^\Omega$ is the $L^2$ Faddeev--Popov metric
adapted to twisted bundles.

Item (iii) is the trivial-sector analogue of
\cite[Thm.~6.1]{RemondiereKRFPB_2026}, with the constant-mode obstruction
removed.
\end{remark}

% =========================================================================
\section{Bakry--\'Emery on the twisted orbit space}\label{sec:BE-twisted}

\subsection{Twisted Faddeev--Popov spectral bound (KR--FP--3 analogue)}

\begin{theorem}[Twisted KR--FP--3 (PROVED UNCONDITIONAL modulo
Hypothesis~\ref{hyp:twist-rigidity})]\label{thm:KRFP3-twisted}
On the twisted lattice $(T^4_L)^\Omega$ with non-trivial twist, every
$A$ in the twisted action slice $\overline\Lambda^\Omega_{S_0}$
satisfies
$$
\lambda_{\min}(M^\Omega[A]) \geq m_\Omega^2 \cdot (1 - \kFP),
\qquad \kFP = 1/(2|\Phi^+(SU(N))|),
$$
\emph{unconditionally on the (H1) generic-vanishing hypothesis} of
\cite{RemondiereKRFP3_2026}, the proof reducing to the spectral structure
of the twisted Hodge Laplacian (Proposition~\ref{prop:Hodge-twisted})
and the Kostant identity (Lemma~KR--FP--2).
\end{theorem}
\begin{proof}[Proof sketch]
The proof of \cite[Thm.~3.1]{RemondiereKRFP3_2026} relies on (H1) only
through the requirement that minimising eigenfunctions of $M[A]$ on the
slice $\overline\Lambda_{S_0}$ are generic in the Cartan-component
sense. On the twisted bundle, the absence of the constant zero
eigenvalue (Proposition~\ref{prop:Hodge-twisted}) means the minimising
eigenfunctions cannot reside in the constant zero subspace; they must
lie in the orthogonal complement, which is exactly the generic
Cartan-component subspace by Lemma~KR--FP--2. Hence (H1) is
\emph{automatic}.

The remaining (H2) compact-manifold Sobolev constant and (H3)
measurable Cartan selection are technical and well-controlled
by the standard arguments of \cite[\S 5]{RemondiereKRFP3_2026}; they
are not affected by the twist.
\end{proof}

\begin{corollary}[Twisted Ricci bound]\label{cor:twisted-Ric}
Under Hypothesis~\ref{hyp:twist-rigidity}(ii) and Theorem~\ref{thm:KRFP3-twisted},
$\Ric_{\mathcal A^\Omega/\mathcal G^\Omega}\restriction_{\overline\Lambda^\Omega_{S_0}}
\geq (1-\kFP) m_\Omega^2 \cdot g^\Omega$.
\end{corollary}

\subsection{Bakry--\'Emery on the twisted manifold}

\begin{theorem}[Twisted KR--FP--B (PROVED UNCONDITIONAL modulo
Hypothesis~\ref{hyp:twist-rigidity})]\label{thm:KRFPB-twisted}
The heat semigroup on the twisted orbit space $\mathcal A^\Omega/
\mathcal G^\Omega$ restricted to $\overline\Lambda^\Omega_{S_0}$
satisfies $\mathrm{CD}((1-\kFP) m_\Omega^2, \infty)$, hence a uniform
LSI with constant $\rho^\Omega \geq (1-\kFP) m_\Omega^2$.

Combined with the twisted-lattice Wilson measure and the trivial-sector
analogue of the BBD continuum descent (which is automatic on the
twisted measure because the constant zero mode is absent), the twisted
Wilson measure $\mu^\Omega_{a,\beta,L}$ satisfies LSI with constant
$$
C_\mathrm{LSI}(\mu^\Omega_{a,\beta,L}) \;\leq\; \cinf(D) / [(1-\kFP) m_\Omega^2].
$$
\end{theorem}
\begin{proof}[Proof sketch]
Bakry--\'Emery on the smooth twisted manifold
$\mathcal A^\Omega/\mathcal G^\Omega$ with $\Ric \geq K\,g$
(Corollary~\ref{cor:twisted-Ric}) gives $\mathrm{CD}(K,\infty)$
\cite[Thm.~4.1]{RemondiereKRFPB_2026}, hence LSI by classical
Bakry--\'Emery~\cite{BakryEmery1985, BakryGentilLedoux2014}. The
trivial-sector continuum descent of~\cite[\S 5]{RemondiereKRFPB_2026}
goes through unchanged on the twisted measure because the constant-mode
obstruction (which was the only one requiring BBD input) is absent.
\end{proof}

% =========================================================================
\section{Main pre-validation theorem}\label{sec:main}

\begin{theorem}[Main pre-validation theorem]\label{thm:main}
Let $G = SU(N)$, $D = 4$, and consider the twisted lattice
$(T^4_L)^\Omega$ with non-trivial twist tensor $n^{\mu\nu} \neq 0
\pmod N$ as in Definition~\ref{def:twist}. Let
$\mu^\Omega_{a,\beta,L}$ be the twisted Wilson measure
\eqref{eq:twisted-Wilson-action}.

Under Hypothesis~\ref{hyp:twist-rigidity}, $\mu^\Omega_{a,\beta,L}$
satisfies, uniformly in $L$ and $a$,
$$
C_\mathrm{LSI}(\mu^\Omega_{a,\beta,L}) \;\leq\; \frac{\cinf(D)}{(1-\kFP) m_\Omega^2}
\;\leq\; \frac{N^2\,\cinf(D)}{(1-\kFP)\,m_0^2}.
$$
For $G = SU(3)$, $D = 4$, $\kFP = 1/6$, $\cinf(4) = 4.05$, this reads
$$
C_\mathrm{LSI}(\mu^\Omega_{a,\beta,L}) \;\leq\; \frac{9 \cdot 4.05}{(5/6)\,m_0^2}
\;\approx\; \frac{43.7}{m_0^2}.
$$
\end{theorem}

\begin{remark}[Comparison with conditional KR--FP--B]\label{rem:compare-KRFP-B}
The constant in Theorem~\ref{thm:main} is a factor of $N^2$ worse than
the corresponding conditional bound of~\cite[Thm.~6.1]{RemondiereKRFPB_2026}
in the trivial sector, but it is \emph{unconditional} (modulo the named
Hypothesis~\ref{hyp:twist-rigidity}). The $N^2$ degradation reflects the
twisted-momentum lattice spacing $2\pi/(NL)$ instead of $2\pi/L$.

For $G=SU(3)$, the twisted bound is a factor $9$ worse but still
$L$-uniform. As $L\to\infty$ both bounds are positive and $O(1/L^2)$
divergent, but on the unit-volume rescaled lattice the bounds give
positive constants, hence positive mass gap in the continuum limit.
\end{remark}

% =========================================================================
\section{Lattice-JAX specification for numerical validation}\label{sec:jax}

\subsection{Recommendation}
We recommend the following lattice-JAX run as the highest-priority
numerical test of Hypothesis~\ref{hyp:twist-rigidity} and the
twisted Wilson measure LSI:

\subsection{Configuration}
\begin{itemize}[leftmargin=2em]
\item Gauge group: $G = SU(3)$.
\item Lattice dimension: $D = 4$.
\item Volumes: $L \in \{8, 12, 16\}$.
\item Couplings: $\beta \in \{2.5, 3.0, 3.5\}$ (continuum regime).
\item Twist matrices: standard 't Hooft pair on directions $(1,2)$:
\begin{itemize}
\item $\Omega_1 = \mathrm{diag}(1, \omega, \omega^2)$, $\omega = e^{2\pi
i/3}$.
\item $\Omega_2 = $ cyclic shift matrix $\mathrm{shift}(3)$.
\item $\Omega_3 = \Omega_4 = \mathrm{Id}$ (twist only in the $(1,2)$ plane).
\end{itemize}
The resulting twist tensor: $n^{12} = 1 \pmod 3$, all others $0$.
\end{itemize}

\subsection{Observables}
\begin{enumerate}[label=(O\arabic*),leftmargin=2em]
\item \textbf{Smallest Faddeev--Popov eigenvalue} $\lambda_{\min}(M^\Omega[A])$
on configurations sampled from $\mu^\Omega_{a,\beta,L}$. Measure as
function of $L$ and $\beta$, verify lower bound $(5/6)\,m_\Omega^2$
with $m_\Omega^2 = (2\pi/(NL))^2$.
\item \textbf{Spectral gap of the Hessian of the twisted Wilson action}
$\Hess(\beta S_W^\Omega)$ at configurations sampled from
$\mu^\Omega_{a,\beta,L}$. Verify constant lower bound uniform in $L$.
\item \textbf{LSI constant of the twisted Wilson measure} estimated via
the Rothaus--Simon mean-field method (sample-based estimate of the
entropy-Dirichlet ratio). Verify $C_\mathrm{LSI}(\mu^\Omega_{a,\beta,L})
\leq 43.7 / m_0^2$ for $L \in \{8,12,16\}$.
\item \textbf{Comparison with periodic measure}: same observables
sampled from $\mu^\mathrm{per}_{a,\beta,L}$ for confirmation that the
twisted measure removes the constant zero mode and the periodic measure
does not.
\end{enumerate}

\subsection{Algorithm}
\begin{enumerate}[label=\arabic*.,leftmargin=2em]
\item Initialise random link configuration $U \in \mathcal C^\Omega$
satisfying~\eqref{eq:twisted-BC}.
\item Thermalise via heat-bath or Metropolis updates with the twisted
Wilson action~\eqref{eq:twisted-Wilson-action}, ensuring
$\langle P\rangle$ matches literature (sanity check).
\item Sample $N_\mathrm{config} = 200$ thermalised configurations.
\item For each configuration:
\begin{enumerate}[label=(\alph*),leftmargin=2em]
\item Compute the FP operator $M^\Omega[A] = d_A^\dagger d_A$
restricted to the twisted bundle.
\item Diagonalise (sparse Krylov) and extract $\lambda_{\min}$.
\item Compute the Hessian of $S_W^\Omega$ and extract its smallest
non-zero eigenvalue.
\end{enumerate}
\item Aggregate over configurations and extract finite-volume
extrapolations.
\end{enumerate}

\subsection{Anticipated outcomes}
\begin{itemize}[leftmargin=2em]
\item For \emph{twisted} measure: $\lambda_{\min}(M^\Omega) \geq (5/6)
m_\Omega^2$ confirmed within $5\%$. Hessian smallest eigenvalue
positive uniformly in $L$. LSI constant bounded as predicted.
\item For \emph{periodic} measure (comparison): $\lambda_{\min}(M) \approx
0$ from the constant zero modes (Pilier 3 sub-3 obstruction confirmed
numerically), Hessian smallest eigenvalue $\to 0$ as $L\to\infty$.
\end{itemize}

\subsection{Code outline}
\begin{verbatim}
import jax
import jax.numpy as jnp

N = 3                      # SU(3)
L = 12                     # lattice side (8, 12, 16)
beta = 3.0                 # coupling
D = 4                      # dimension

# Standard 't Hooft pair
omega = jnp.exp(2j * jnp.pi / N)
Omega1 = jnp.diag(jnp.array([1, omega, omega**2]))
Omega2 = jnp.roll(jnp.eye(N), 1, axis=0)
Omega3 = jnp.eye(N)
Omega4 = jnp.eye(N)

# Twisted plaquette phase
z12 = jnp.exp(-2j * jnp.pi / N)   # twist tensor n^{12} = 1 mod N

def twisted_plaquette(U, x, mu, nu, L, Omegas):
    # standard plaquette holonomy, with twist phase z_{mu nu}
    # at links crossing the boundary
    Umu  = U[x[0], x[1], x[2], x[3], mu]
    Unu  = U[(x[0]+(mu==0))%L, (x[1]+(mu==1))%L,
             (x[2]+(mu==2))%L, (x[3]+(mu==3))%L, nu]
    # ... (full implementation requires careful index management
    # for the twist phase at boundary plaquettes)
    return jnp.trace(Umu @ Unu @ Umu_dag @ Unu_dag) * z12 if boundary else 1

def twisted_wilson_action(U, beta, N, L, twist_pair=(0,1)):
    # sum over plaquettes of 1 - Re Tr(z * plaq)/N
    # ...
    pass

def faddeev_popov_M(U):
    # bundle-twisted FP operator M = d_A^dag d_A
    # acting on adjoint-valued 0-forms with twisted BC
    # ...
    pass

# Sample thermalised configurations and measure smallest eigenvalue
def main():
    key = jax.random.PRNGKey(42)
    U = init_random_twisted(key, L, N, Omegas=[Omega1, Omega2, Omega3, Omega4])
    U = thermalise(U, beta, n_steps=5000)
    for cfg in range(N_config := 200):
        U = update(U, beta, n_sweeps=10)
        M = faddeev_popov_M(U)
        lam_min = jnp.linalg.eigvalsh(M)[0]   # smallest eigenvalue
        # ... aggregate and report
\end{verbatim}

\subsection{Cost estimate}
\begin{itemize}[leftmargin=2em]
\item Lattice volume $L = 12$, $N = 3$, $D = 4$: $4 \cdot 12^4 \cdot 8 = 663{,}552$
link real degrees of freedom.
\item FP operator dimension: $(N^2-1)\cdot L^4 = 8 \cdot 20736 = 165{,}888$.
\item Sparse Krylov: feasible on consumer GPU (e.g.\ RTX 3090, $\sim 1$~min
per configuration for $\lambda_{\min}$).
\item Total: $200$ configurations $\times \sim 1$~min $= \sim 3$~hours
per $(L,\beta)$ tuple. Total parameter sweep $(3 \times 3 = 9)$:
$\sim 27$ hours $\sim 1$~day on a Vast.AI RTX 3090 ($\sim \$5$).
\item Implementation: $\sim 1$~week of dev/test, $1$~day of run, $\sim 2$~days
of analysis.
\end{itemize}

\subsection{Priority recommendation}
We recommend this lattice run as the \emph{highest-priority short-term
test} ($2-4$ months ETA) of the trivial-sector obstruction-removal
mechanism. If the test confirms Hypothesis~\ref{hyp:twist-rigidity},
Theorem~\ref{thm:main} is unconditional, and the structural
contribution of this paper is publishable standalone as a clean
removal-of-the-constant-zero-mode result for a measurable subset of the
configuration space (the twisted bundle).

% =========================================================================
\section{Limits, outlook, and recommendation}\label{sec:outlook}

\subsection{What is proved unconditionally}
\begin{itemize}[leftmargin=2em]
\item Lemma~\ref{lem:centraliser}: centraliser computation.
\item Corollary~\ref{cor:twist-eats-constant}: twist eats the constant
gauge mode.
\item Proposition~\ref{prop:Hodge-twisted}: twisted Hodge Laplacian
has no zero eigenvalue, smallest eigenvalue $\geq (2\pi/(NL))^2$.
\end{itemize}

\subsection{What is conditional on Hypothesis~\ref{hyp:twist-rigidity}}
\begin{itemize}[leftmargin=2em]
\item Theorem~\ref{thm:KRFP3-twisted}: twisted-KR--FP--3 spectral
bound on $\overline\Lambda^\Omega_{S_0}$.
\item Theorem~\ref{thm:KRFPB-twisted}: twisted-KR--FP--B Bakry--\'Emery
LSI bound.
\item Theorem~\ref{thm:main}: main pre-validation theorem.
\end{itemize}

\subsection{Why Hypothesis~\ref{hyp:twist-rigidity} is plausible}
\begin{itemize}[leftmargin=2em]
\item Item (i): the structural reason for the zero mode is
\emph{exclusively} the constant gauge transformation, which is killed
by the twist (Corollary~\ref{cor:twist-eats-constant}). The other
analytical inputs to KR--FP--3 (Birman--Schwinger, Aubin--Talenti,
Kostant identity) are bundle-independent.
\item Item (ii): Babelon--Viallet O'Neill is a Riemannian-geometric
identity on the principal bundle $\mathcal A \to \mathcal A/\mathcal G$;
it goes through identically on twisted bundles by replacing the trivial
$P$ with the twisted $P^\Omega$.
\item Item (iii): the Bakry--\'Emery $\Gamma_2$-calculus is a
purely Riemannian-geometric framework on the orbit space, and on the
twisted manifold (where the gauge orbits are now twisted-bundle orbits)
the calculation goes through unchanged.
\end{itemize}

\subsection{Why a clean rigorous verification is missing}
The reason this is currently a \emph{hypothesis} rather than a
\emph{theorem} is that the existing literature on 't Hooft twist
focuses on \emph{computational} aspects (twisted Eguchi--Kawai
\cite{EguchiKawai1982}, large-$N$ reduction
\cite{GonzalezArroyoOkawa1983}, citation to verify) rather than on the
\emph{analytical} (Bakry--\'Emery, LSI) consequences for the Wilson
measure. A clean rigorous derivation of items (i)--(iii) above
requires combining the geometric construction of twisted bundles
(van Baal~\cite{vanBaal1996}, citation to verify) with the
Bakry--\'Emery framework of~\cite{RemondiereKRFPB_2026} --- a
combination not present in the existing literature but well within
reach of contemporary techniques.

\subsection{Why the lattice run is the recommended path}
A direct numerical verification of items (i)--(iii) on
$G = SU(3), L \in \{8,12,16\}$ takes $\sim 1$ day on a GPU and
costs $\sim \$5$. If the test confirms the hypothesis, the structural
contribution of this paper is publishable as a clean result. If the
test \emph{contradicts} the hypothesis (e.g.\ a subtle twist-induced
analytical pathology), we identify a new open problem to be solved.

\subsection{Combined with the AHS instanton paper, the structural picture}
\begin{itemize}[leftmargin=2em]
\item Companion paper \emph{AHS-Instanton-LSI}~\cite{RemondiereAHSInstantonLSI_2026}
gives unconditional LSI on \emph{non-trivial} topological sectors
$k\neq 0$ of the \emph{periodic} Wilson measure.
\item This paper gives conditional (on twist-rigidity hypothesis)
unconditional-on-(H1) LSI on the \emph{trivial} topological sector $k=0$
of the \emph{twisted} Wilson measure.
\item Together, these two structural results cover the full
configuration space \emph{partitioned} into (i) periodic + non-trivial
topology and (ii) twisted + trivial topology. The periodic + trivial
sector remains conditional on (H1)--(H3) of
\cite{RemondiereKRFP3_2026} plus BBD~LSI.
\end{itemize}

\subsection{Honest probability estimates}
\begin{itemize}[leftmargin=2em]
\item P(Hypothesis~\ref{hyp:twist-rigidity} valid): $80-90\%$ on
structural grounds + (i)--(iii) above.
\item P(lattice JAX run confirms Hypothesis): $80-90\%$ if hypothesis
is valid, else $\leq 30\%$.
\item P(rigorous mathematical proof of Hypothesis): $50-70\%$ on the
$1-3$y horizon (with adapted Babelon--Viallet + Bakry--\'Emery on
twisted bundles).
\item P(paper publishable in LMP/CMP): $70-85\%$ as a structural
pre-validation + lattice JAX specification.
\end{itemize}

\subsection{Recommendation}
We recommend:
\begin{enumerate}[label=\arabic*.,leftmargin=2em]
\item Launch the lattice-JAX run \emph{this week} ($1$ day implementation
+ $1$ day run + $1$ day analysis).
\item If the run confirms the hypothesis, finalise the paper and
submit to LMP.
\item In parallel, pursue the rigorous mathematical proof of
Hypothesis~\ref{hyp:twist-rigidity} (1-3 month effort).
\item Coordinate with the AHS-instanton paper for joint publication
or sequential submission as a series of two structural notes.
\end{enumerate}

% =========================================================================
\section*{Acknowledgments}
This work was prepared with computational research assistance from
Anthropic's Claude system (Opus~4.7, 1M context window). All
mathematical content (theorems, proofs, hypotheses, references) is
the sole responsibility of the author.

% =========================================================================
\begin{thebibliography}{99}

\bibitem{BakryEmery1985}
D.~Bakry, M.~\'Emery,
``Diffusions hypercontractives,''
\emph{S\'eminaire de Probabilit\'es XIX},
Lecture Notes in Mathematics \textbf{1123} (Springer, 1985) 177--206.

\bibitem{BakryGentilLedoux2014}
D.~Bakry, I.~Gentil, M.~Ledoux,
\emph{Analysis and Geometry of Markov Diffusion Operators},
Grundlehren \textbf{348}, Springer, 2014.

\bibitem{BauerschmidtDagallier2024}
R.~Bauerschmidt, B.~Dagallier,
``Log-Sobolev inequality for the $\varphi^4_2$ and $\varphi^4_3$ measures,''
arXiv:2202.02295.

\bibitem{EguchiKawai1982}
T.~Eguchi, H.~Kawai,
``Reduction of dynamical degrees of freedom in the large-$N$ gauge theory,''
\emph{Phys. Rev. Lett.} \textbf{48} (1982) 1063--1066.

\bibitem{GonzalezArroyoOkawa1983}
A.~Gonz\'alez-Arroyo, M.~Okawa,
``A twisted model for large-$N$ lattice gauge theory,''
\emph{Phys. Lett. B} \textbf{120} (1983) 174--178.
\emph{Citation to verify}: standard textbook reference, exact pages
require direct check.

\bibitem{OpusPolchinskiSubgaps2026}
K.~R\'emondi\`ere,
\emph{Polchinski sub-gaps for SU(N) extension: 5-sub-gap attack with
honest verdicts} (internal note
\texttt{OPUS2\_POLCHINSKI\_SUBGAPS\_2026-05-26}), May 2026.

\bibitem{RemondiereAHSInstantonLSI_2026}
K.~R\'emondi\`ere,
\emph{Logarithmic Sobolev inequality for the instanton sector of pure
SU(N) Wilson lattice gauge theory}
(companion preprint
\texttt{Paper\_AHS\_Instanton\_LSI\_CMP}), May 2026.

\bibitem{RemondiereKRFP3_2026}
K.~R\'emondi\`ere,
\emph{Spectral bound for the Faddeev--Popov operator on the fundamental
modular domain of the slice $\overline\Lambda_{S_0}$: a Birman--Schwinger /
Aubin--Talenti analysis} (companion preprint
\texttt{Paper\_KR\_FP3\_AnnalsMath}), May 2026.

\bibitem{RemondiereKRFPB_2026}
K.~R\'emondi\`ere,
\emph{Mass gap for 4D pure Yang--Mills via Bakry--\'Emery on the
fundamental modular domain: a conditional reduction} (companion preprint
\texttt{Paper\_KR\_FP\_B\_BakryEmery\_LMP}), May 2026.

\bibitem{RemondiereWilsonFlowVoieB_2026}
K.~R\'emondi\`ere,
\emph{Wilson flow trivialisation gives a uniform logarithmic Sobolev
inequality for $SU(N)$ lattice Yang--Mills: a conditional structural
reduction} (companion preprint
\texttt{Paper\_WilsonFlow\_VoieB\_LMP}), May 2026.

\bibitem{tHooft1979}
G.~'t Hooft,
``A property of electric and magnetic flux in non-abelian gauge theories,''
\emph{Nucl. Phys. B} \textbf{153} (1979) 141--160.

\bibitem{tHooft1981}
G.~'t Hooft,
``Some twisted self-dual solutions for the Yang--Mills equations on a
hypertorus,''
\emph{Commun. Math. Phys.} \textbf{81} (1981) 267--275.

\bibitem{vanBaal1996}
P.~van Baal,
``Instantons versus monopoles,''
\emph{Nucl. Phys. B Proc. Suppl.} \textbf{47} (1996) 326--332.
\emph{Citation to verify}: van Baal's work on calorons and twisted
boundary conditions through the 1990s; specific arXiv ID requires
direct verification.

\end{thebibliography}

\end{document}
